\documentclass[10pt,twocolumn,a4paper]{article}

% Page layout
\usepackage[margin=1.8cm, columnsep=0.6cm]{geometry}

% Essential packages
\usepackage{amsmath,amssymb,amsthm}
\usepackage{mathtools}
\usepackage{bm}
\usepackage{graphicx}
\graphicspath{{figures/}}
\usepackage{hyperref}
\usepackage{natbib}
\usepackage{physics}
\usepackage{booktabs}
\usepackage{array}
\usepackage{multirow}
\usepackage{enumitem}
\usepackage{xcolor}
\usepackage{float}
\usepackage{algorithm}
\usepackage{algorithmic}
\usepackage{tikz}
\usepackage{tikz-cd}
\usetikzlibrary{arrows,positioning,calc}

% Font and typography
\usepackage[utf8]{inputenc}
\usepackage[T1]{fontenc}
\usepackage{lmodern}
\usepackage{microtype}

% Theorem environments
\theoremstyle{plain}
\newtheorem{theorem}{Theorem}[section]
\newtheorem{lemma}[theorem]{Lemma}
\newtheorem{proposition}[theorem]{Proposition}
\newtheorem{corollary}[theorem]{Corollary}

\theoremstyle{definition}
\newtheorem{definition}[theorem]{Definition}
\newtheorem{axiom}{Axiom}
\newtheorem{assumption}[theorem]{Assumption}

\theoremstyle{remark}
\newtheorem{remark}[theorem]{Remark}
\newtheorem{example}[theorem]{Example}

% Hyperref setup
\hypersetup{
    colorlinks=true,
    linkcolor=blue,
    citecolor=blue,
    urlcolor=blue,
    pdftitle={Sequential Constraint Propagation in Fuzzy Cellular Circuits},
    pdfauthor={Kundai Farai Sachikonye}
}

% Custom commands
\newcommand{\kB}{k_{\mathrm{B}}}
\newcommand{\Depth}{\mathcal{M}}
\newcommand{\fuzzy}{\widetilde}
\newcommand{\node}{\mathsf{v}}
\newcommand{\edge}{\mathsf{e}}
\newcommand{\circuit}{\mathcal{G}}
\newcommand{\resolve}{\mathsf{R}}
\newcommand{\valid}{\mathsf{V}}
\newcommand{\consist}{\mathsf{C}}
\newcommand{\sig}{\sigma}
\newcommand{\fuzzyval}{\mu}
\newcommand{\membership}{\mu}
\newcommand{\tpart}{\tau_p}
\newcommand{\order}{\langle r \rangle}
\newcommand{\R}{\mathbb{R}}
\newcommand{\N}{\mathbb{N}}
\newcommand{\Z}{\mathbb{Z}}
\newcommand{\dcat}{d_{\mathrm{C}}}
\newcommand{\Svec}{\mathbf{S}}
\newcommand{\defect}{\mathcal{D}}
\newcommand{\window}{\mathcal{W}}
\newcommand{\topology}{\mathcal{T}}
\newcommand{\landscape}{\mathcal{L}}
\newcommand{\prop}{\rightsquigarrow}

\begin{document}

% Title
\title{\textbf{Sequential Constraint Propagation in Fuzzy Cellular Circuits:\\[0.5em]
\large A Step-Discrete Framework for Disease as Topological Inconsistency}}

\author{
Kundai Farai Sachikonye\\
Technical University of Munich, School of Life Sciences\\
\texttt{kundai.sachikonye@wzw.tum.de}
}

\date{\today}

\maketitle

%==============================================================================
\begin{abstract}
%==============================================================================
\noindent We present a framework for cellular dynamics based on sequential constraint propagation through fuzzy-valued circuit networks. The central observation is that cells do not maintain a reference state against which to detect disease---they operate by local, step-wise signal validation where each resolved signal constrains the resolution of the next. We model the cell as a circuit graph $\circuit = (V, E)$ where each node $\node_i$ carries a fuzzy membership value $\fuzzyval_i \in [0,1]$ representing its undetermined state. Nodes are resolved sequentially: solving node $\node_i$ collapses its fuzzy value to a discrete output, which propagates as a boundary condition to adjacent nodes. We prove three principal results. First, that step-discrete sequential resolution is the only dynamics consistent with the absence of a global reference state (Theorem~\ref{thm:no_template}). Second, that disease corresponds to topological inconsistency---closed loops in the circuit graph that fail to return self-consistent values after complete traversal (Theorem~\ref{thm:disease_inconsistency}). Third, that such inconsistencies are locally undetectable: each individual step passes validation, yet the global circuit is inconsistent (Theorem~\ref{thm:local_invisibility}). This explains why cells are ``blind'' to their own disease. We derive a consistency index $\consist(\circuit)$ that quantifies global circuit health from local validation residuals, establish its relationship to measurable macroscopic signals (pH, volume, membrane potential), and show that disease progression corresponds to monotonic decrease of $\consist(\circuit)$ through a sequence of individually valid steps. The framework is validated against glycolytic flux data, mitochondrial dysfunction cascades, and protein misfolding disease progression, reproducing observed dynamics with zero free parameters.
\end{abstract}

%==============================================================================
\section{Introduction}
\label{sec:introduction}
%==============================================================================

\subsection{The Blindness Problem}

A cell contains the complete molecular machinery to execute thousands of simultaneous biochemical reactions. It monitors pH, osmolarity, membrane potential, metabolite concentrations, and gene expression in real time. Yet cells become diseased. A neuron accumulating misfolded protein aggregates continues to function---passing local quality checks at every step---until catastrophic failure. A cancer cell passes every internal checkpoint while violating the constraints of the organism.

This presents a paradox. If cells have access to all information about their internal state, why are they blind to disease?

The standard explanation invokes specific molecular failures: a mutated tumour suppressor, a misfolded chaperone, a broken signalling cascade. But this merely pushes the question back. \emph{Why} can't the cell detect these failures, given that it has complete access to the molecular state?

We propose that the answer is structural. Cells do not operate by comparing their state against a stored template of ``health.'' They operate by \emph{sequential constraint propagation}: each biochemical event validates against the output of the previous event, not against a global reference. Disease is invisible because it accumulates through a sequence of individually valid steps.

\subsection{Why Simulation Fails}

The dominant paradigm in computational biology is simulation: construct a model of cellular reactions and integrate the equations of motion forward in time \citep{karr2012,kitano2002}. The implicit assumption is that a sufficiently detailed simulation will capture disease.

This assumption is flawed for a fundamental reason. A perfect simulation of every reaction in a cell would reproduce the cell's behaviour exactly---\emph{including its blindness to disease}. The simulation would show every reaction proceeding normally, every checkpoint passing, every local signal within bounds, while the global state drifts toward pathology. The disease is not in any individual reaction. It is in the \emph{relational structure} between reactions---a structure that neither the cell nor the simulation explicitly represents.

This observation motivates a different approach. Rather than simulating reactions in continuous time, we model the cell as a \emph{circuit}---a network of nodes connected by constraint relationships---and study how information propagates through this circuit in discrete steps.

\subsection{The Circuit Analogy}

Consider an electronic circuit. Each component (resistor, capacitor, transistor) has a local constitutive relation. The state of the circuit is determined not by simulating electron flow but by solving constraints: Kirchhoff's laws at each node, constitutive relations at each component. The solution propagates through the circuit---solving one node provides boundary conditions for the next.

We propose that cellular biochemistry operates analogously. Each reaction node has a local ``constitutive relation''---the relationship between its inputs and outputs. The cellular state is determined not by simulating mass transfer but by propagating constraints through the reaction network. The critical difference from electronic circuits is that biological nodes carry \emph{fuzzy} values: their state is not precisely determined until resolved against neighbouring constraints.

\subsection{Contributions}

This paper makes the following contributions:
\begin{enumerate}[label=(\roman*), leftmargin=*]
    \item A formal model of cellular dynamics as sequential constraint propagation through fuzzy-valued circuit graphs (Section~\ref{sec:circuit_model})
    \item A proof that step-discrete resolution is the unique dynamics consistent with absence of a global reference state (Section~\ref{sec:dynamics})
    \item A characterization of disease as topological inconsistency in the circuit graph (Section~\ref{sec:disease})
    \item A proof that such inconsistencies are locally undetectable, explaining cellular blindness to disease (Section~\ref{sec:invisibility})
    \item A consistency index derived from measurable macroscopic signals (Section~\ref{sec:consistency})
    \item Validation against glycolytic flux, mitochondrial dysfunction, and protein misfolding data (Section~\ref{sec:validation})
\end{enumerate}

\subsection{Paper Organization}

Part~I (Sections~\ref{sec:foundations}--\ref{sec:dynamics}) establishes the mathematical foundations: bounded phase space, fuzzy circuit graphs, and sequential resolution dynamics. Part~II (Sections~\ref{sec:disease}--\ref{sec:invisibility}) develops the theory of disease as topological inconsistency and proves local undetectability. Part~III (Sections~\ref{sec:consistency}--\ref{sec:therapeutic}) derives the consistency index and its relationship to macroscopic observables. Part~IV (Sections~\ref{sec:validation}--\ref{sec:predictions}) presents validation and experimental predictions. Section~\ref{sec:conclusion} concludes.


%==============================================================================
\part*{I. Foundations}
%==============================================================================

%==============================================================================
\section{First Principles}
\label{sec:foundations}
%==============================================================================

\subsection{Bounded Phase Space}

We begin from a single physical observation.

\begin{axiom}[Boundedness]
\label{ax:bounded}
Every localized physical system occupies a bounded region of phase space:
\begin{equation}
    \Omega \subset \R^{2n}, \quad \mathrm{Vol}(\Omega) < \infty
\end{equation}
where $n$ is the number of degrees of freedom.
\end{axiom}

Cells have finite volume, finite numbers of molecules, and finite energy. This boundedness is not a modelling assumption---it is a physical fact. A bounded system has a finite number of distinguishable states:
\begin{equation}
    N_{\max} = \frac{\mathrm{Vol}(\Omega)}{h^n} < \infty
\end{equation}
where $h^n$ is the minimum phase space cell volume from the uncertainty principle \citep{landau1980}.

\subsection{Partition Structure}

A bounded phase space admits decomposition into distinguishable subregions.

\begin{axiom}[Finite Partition]
\label{ax:partition}
A bounded phase space $\Omega$ admits partition into a finite number of distinguishable cells:
\begin{equation}
    \Omega = \bigsqcup_{i=1}^{N} \Omega_i, \quad N \leq N_{\max}
\end{equation}
\end{axiom}

The number of levels of hierarchical partitioning defines the \emph{partition depth}:

\begin{definition}[Partition Depth]
\label{def:depth}
The partition depth $\Depth$ of a state is the number of hierarchical distinctions required for unique specification:
\begin{equation}
    \Depth = \sum_{i} \log_2(k_i)
\end{equation}
where $k_i$ is the branching factor at level $i$ of the partition hierarchy.
\end{definition}

Partition depth connects directly to thermodynamic entropy:
\begin{equation}
    S = \kB \ln 2 \cdot \Depth
    \label{eq:depth_entropy}
\end{equation}

\subsection{The Absence of Templates}

We now establish a crucial negative result.

\begin{theorem}[No Template Theorem]
\label{thm:no_template}
A bounded physical system cannot store an internal representation of its own complete state.
\end{theorem}

\begin{proof}
Let $\Omega$ be the phase space of the system with $N_{\max}$ distinguishable states. A complete state description requires specifying which of the $N_{\max}$ states the system occupies, requiring $\log_2(N_{\max})$ bits. Storing this description within the system requires dedicating a subsystem with at least $\log_2(N_{\max})$ distinguishable states to the representation.

But this storage subsystem is itself part of the system. Its state must also be represented, requiring additional storage of $\log_2(\log_2(N_{\max}))$ bits, and so on. The total storage requirement diverges:
\begin{equation}
    S_{\mathrm{total}} = \log_2(N_{\max}) + \log_2(\log_2(N_{\max})) + \cdots
\end{equation}

Since the system has finite capacity $N_{\max}$, it cannot store a complete self-representation. Any internal model is necessarily incomplete.
\end{proof}

\begin{corollary}[No Health Template]
\label{cor:no_health}
A cell cannot store a complete representation of its ``healthy state.'' Any comparison-based disease detection mechanism is necessarily incomplete.
\end{corollary}

This corollary is the formal statement of the blindness problem. The cell cannot detect disease by comparing its current state to a stored reference, because no such reference can exist within the cell.


%==============================================================================
\section{Fuzzy Circuit Graphs}
\label{sec:circuit_model}
%==============================================================================

\subsection{The Circuit Graph}

Having established that cells cannot operate by state comparison, we develop the framework for how they \emph{do} operate.

\begin{definition}[Cellular Circuit Graph]
\label{def:circuit}
A cellular circuit graph is a directed graph $\circuit = (V, E, \fuzzyval, \Phi)$ where:
\begin{enumerate}[label=(\roman*)]
    \item $V = \{\node_1, \node_2, \ldots, \node_n\}$ is the set of reaction nodes
    \item $E \subseteq V \times V$ is the set of constraint edges
    \item $\fuzzyval: V \to [0, 1]$ assigns a fuzzy membership value to each node
    \item $\Phi: E \to \mathcal{F}$ assigns a constraint function to each edge, where $\mathcal{F}$ is the space of functions $[0,1] \to [0,1]$
\end{enumerate}
\end{definition}

Each node represents a reaction or cluster of reactions. The fuzzy value $\fuzzyval_i \in [0,1]$ represents the degree to which the node's state is determined: $\fuzzyval_i = 0$ means completely undetermined, $\fuzzyval_i = 1$ means fully resolved to a definite value.

\begin{definition}[Node State]
\label{def:node_state}
The state of node $\node_i$ is a pair $(\fuzzyval_i, x_i)$ where:
\begin{itemize}
    \item $\fuzzyval_i \in [0,1]$ is the membership value (degree of determination)
    \item $x_i \in \R^{d_i}$ is the discrete value, defined only when $\fuzzyval_i = 1$
\end{itemize}
When $\fuzzyval_i < 1$, $x_i$ is a fuzzy set on $\R^{d_i}$ with membership function $\mu_{x_i}$.
\end{definition}

\begin{remark}
The dimension $d_i$ depends on the node. For a simple metabolic reaction, $d_i$ might be 1 (a flux value). For a signalling hub, $d_i$ might be large, encoding multiple output channels.
\end{remark}

\subsection{Fuzzy Values and Biological Observables}

The fuzzy membership value $\fuzzyval_i$ corresponds to a physical quantity: the degree to which the reaction at node $\node_i$ is constrained by its context.

\begin{definition}[Macroscopic Signals]
\label{def:signals}
A macroscopic signal $\sig_k$ is a coarse-grained observable that is cheaply measurable from the node's local environment:
\begin{equation}
    \sig_k: V \to \R, \quad k \in \{\text{pH}, \text{vol}, \Delta\psi, \text{hydro}, \ldots\}
\end{equation}
We denote the signal vector at node $\node_i$ as $\bm{\sig}_i = (\sig_1(\node_i), \sig_2(\node_i), \ldots, \sig_K(\node_i))$.
\end{definition}

These signals are the ``clean'' observables: pH, volume changes, membrane potential $\Delta\psi$, hydrophobicity gradients---macroscopic quantities that integrate over many molecular degrees of freedom and are therefore robust to microscopic fluctuations.

\begin{theorem}[Signal Sufficiency]
\label{thm:signal_sufficiency}
For sequential constraint propagation, macroscopic signals are sufficient. The fuzzy value of node $\node_i$ is determined by the signal vector of its resolved predecessors:
\begin{equation}
    \fuzzyval_i = f(\bm{\sig}_j : \node_j \in \mathrm{pred}(\node_i), \, \fuzzyval_j = 1)
    \label{eq:fuzzy_from_signals}
\end{equation}
where $\mathrm{pred}(\node_i) = \{\node_j : (\node_j, \node_i) \in E\}$ is the set of predecessor nodes.
\end{theorem}

\begin{proof}
Each node's constitutive relation maps input signals to output signals. When predecessor nodes are resolved ($\fuzzyval_j = 1$), their signal vectors $\bm{\sig}_j$ are definite. The constraint functions $\Phi_{ji}$ on the edges $(\node_j, \node_i)$ propagate these definite values as boundary conditions for node $\node_i$. The fuzzy value of $\node_i$ then depends only on how well these boundary conditions constrain $\node_i$'s state---which is determined entirely by the signal vectors of resolved predecessors.
\end{proof}

\subsection{Constraint Functions}

The edges of the circuit graph carry constraint functions that encode how one node's resolution affects its neighbours.

\begin{definition}[Constraint Function]
\label{def:constraint}
The constraint function $\Phi_{ij}: [0,1] \times \R^{d_i} \to [0,1] \times \R^{d_j}$ on edge $(\node_i, \node_j)$ maps the resolved state of $\node_i$ to a restriction on the state of $\node_j$:
\begin{equation}
    (\fuzzyval_j', x_j') = \Phi_{ij}(\fuzzyval_i, x_i)
\end{equation}
where $\fuzzyval_j' \geq \fuzzyval_j$ (resolution can only increase) and $x_j'$ is the updated fuzzy set for $\node_j$.
\end{definition}

\begin{proposition}[Monotonicity]
\label{prop:monotone}
Constraint propagation is monotonic in fuzzy value:
\begin{equation}
    \fuzzyval_j^{(t+1)} \geq \fuzzyval_j^{(t)}
\end{equation}
Nodes become more determined, never less determined, as constraints propagate.
\end{proposition}

\begin{proof}
Each resolved predecessor provides additional boundary conditions. Additional constraints can only reduce the set of compatible states for $\node_j$, which increases $\fuzzyval_j$ (higher determination). Formally, if $\mathcal{X}_j^{(t)}$ is the set of states compatible with $t$ resolved predecessors, then $\mathcal{X}_j^{(t+1)} \subseteq \mathcal{X}_j^{(t)}$, and $\fuzzyval_j = 1 - |\mathcal{X}_j|/|\mathcal{X}_j^{(0)}|$.
\end{proof}

\subsection{Circuit Topology}

The topology of the circuit graph encodes the dependency structure of cellular reactions.

\begin{definition}[Circuit Topology Classes]
\label{def:topology}
A circuit graph $\circuit$ contains:
\begin{enumerate}[label=(\roman*)]
    \item \textbf{Feed-forward paths:} directed acyclic subgraphs (metabolic pathways)
    \item \textbf{Cycles:} closed directed paths (feedback loops, metabolic cycles)
    \item \textbf{Hubs:} high-degree nodes (signalling centres, ATP/NAD$^+$ pools)
    \item \textbf{Bridges:} edges whose removal disconnects the graph (essential reactions)
\end{enumerate}
\end{definition}

\begin{definition}[Loop]
\label{def:loop}
A loop $\ell$ in $\circuit$ is a closed directed path $\node_{i_1} \to \node_{i_2} \to \cdots \to \node_{i_k} \to \node_{i_1}$. The \emph{loop length} is $|\ell| = k$.
\end{definition}

Loops are critical because they are where consistency can be tested: a signal that propagates around a loop must return a value consistent with its starting value. This is the origin of both self-regulation and disease vulnerability.


%==============================================================================
\section{Sequential Resolution Dynamics}
\label{sec:dynamics}
%==============================================================================

\subsection{The Resolution Operator}

We now define how the circuit evolves.

\begin{definition}[Resolution Operator]
\label{def:resolve}
The resolution operator $\resolve_i$ acts on the circuit state by resolving node $\node_i$:
\begin{equation}
    \resolve_i: (\fuzzyval_i, \{x_i\}) \mapsto (1, x_i^*)
\end{equation}
where $x_i^*$ is the unique discrete value satisfying all constraints from resolved predecessors:
\begin{equation}
    x_i^* = \arg\min_{x \in \mathcal{X}_i} \sum_{j \in \mathrm{pred}(i)} \|\Phi_{ji}(x_j^*) - x\|^2
    \label{eq:resolution}
\end{equation}
\end{definition}

Resolution is a \emph{local} operation: it depends only on the resolved values of predecessor nodes. It requires no global information about the circuit state.

\begin{definition}[Resolution Sequence]
\label{def:sequence}
A resolution sequence is an ordered list of resolution operations:
\begin{equation}
    \mathcal{S} = (\resolve_{i_1}, \resolve_{i_2}, \ldots, \resolve_{i_n})
\end{equation}
The sequence is \emph{valid} if each node is resolved only after all its predecessors (with respect to the current resolution step) have been resolved.
\end{definition}

\begin{theorem}[Sequential Resolution is Necessary]
\label{thm:sequential}
In the absence of a global reference state, the only well-defined dynamics on a fuzzy circuit graph is sequential constraint propagation.
\end{theorem}

\begin{proof}
Consider alternative dynamics:

\textbf{(i) Simultaneous resolution.} Resolving all nodes simultaneously requires a global consistency condition---a system of equations coupling all nodes. Solving this system requires representing the global state, which by Theorem~\ref{thm:no_template} is impossible within the system.

\textbf{(ii) Continuous-time evolution.} Continuous dynamics $\dot{x}_i = F_i(\{x_j\})$ requires evaluating all coupling terms simultaneously at each instant. This is equivalent to simultaneous resolution in the limit $\Delta t \to 0$.

\textbf{(iii) Sequential resolution.} Each resolution step $\resolve_i$ depends only on the finite set of resolved predecessors. No global state representation is required. The dynamics is well-defined at each step.

Since (i) and (ii) require global state representation and (iii) does not, sequential resolution is the unique dynamics consistent with Theorem~\ref{thm:no_template}.
\end{proof}

\subsection{The Partition Lag}

Each resolution step requires a minimum time.

\begin{definition}[Partition Lag]
\label{def:partition_lag}
The partition lag $\tpart$ is the irreducible time required for one resolution operation:
\begin{equation}
    \tpart = \frac{\hbar}{\kB T}
    \label{eq:partition_lag}
\end{equation}
This is the quantum-mechanical minimum time for a state distinction at thermal energy $\kB T$.
\end{definition}

At $T = 310$~K (physiological temperature):
\begin{equation}
    \tpart = \frac{1.055 \times 10^{-34}}{1.38 \times 10^{-23} \times 310} \approx 2.5 \times 10^{-14}~\text{s}
\end{equation}

The effective resolution time for a biological node is much longer, as it includes the physical reorganization of the reaction:
\begin{equation}
    \tau_{\mathrm{eff}} = \tpart + \tau_{\mathrm{reorg}}
    \label{eq:effective_lag}
\end{equation}
where $\tau_{\mathrm{reorg}}$ ranges from $\sim 10^{-12}$~s (fast enzyme conformational changes) to $\sim 10^{0}$~s (slow signalling cascades).

\subsection{Step-Discrete Time}

The circuit does not evolve in continuous time. It evolves in \emph{steps}.

\begin{definition}[Circuit Step]
\label{def:step}
A circuit step $t \to t+1$ consists of:
\begin{enumerate}[label=(\roman*)]
    \item Selection of the next node $\node_{i_{t+1}}$ to resolve
    \item Resolution: $\resolve_{i_{t+1}}$ applied to the current circuit state
    \item Propagation: boundary conditions updated for all successors of $\node_{i_{t+1}}$
\end{enumerate}
The circuit state at step $t$ is:
\begin{equation}
    \circuit^{(t)} = \left(\{\fuzzyval_i^{(t)}\}, \{x_i^{(t)}\}, E, \Phi\right)
\end{equation}
\end{definition}

\begin{theorem}[Discrete Dynamics]
\label{thm:discrete}
The circuit state $\circuit^{(t+1)}$ depends on $\circuit^{(t)}$ and the resolution $\resolve_{i_{t+1}}$ alone:
\begin{equation}
    \circuit^{(t+1)} = \resolve_{i_{t+1}} \circ \circuit^{(t)}
    \label{eq:discrete_dynamics}
\end{equation}
The dynamics is first-order Markov in steps.
\end{theorem}

\begin{proof}
Resolution of $\node_{i_{t+1}}$ depends only on the current resolved values of its predecessors (Equation~\ref{eq:resolution}). These are contained in $\circuit^{(t)}$. No information from earlier steps is required beyond what is already encoded in the current state.
\end{proof}

\subsection{The Validation Function}

At each step, the resolved value is validated against the incoming constraints.

\begin{definition}[Local Validation]
\label{def:validation}
The local validation function $\valid_i$ at node $\node_i$ is:
\begin{equation}
    \valid_i(x_i^*) = \prod_{j \in \mathrm{pred}(i)} \Phi_{ji}^{\mathrm{compat}}(x_j^*, x_i^*)
    \label{eq:validation}
\end{equation}
where $\Phi_{ji}^{\mathrm{compat}}: \R^{d_j} \times \R^{d_i} \to [0,1]$ is the compatibility function on edge $(\node_j, \node_i)$. The node passes validation if $\valid_i \geq \theta$ for threshold $\theta \in (0,1)$.
\end{definition}

\begin{remark}
The product structure in Equation~\ref{eq:validation} means that a node passes validation if and only if it is compatible with \emph{each} of its resolved predecessors. A single incompatibility ($\Phi_{ji}^{\mathrm{compat}} \approx 0$ for some $j$) causes validation failure regardless of all other constraints.
\end{remark}

\begin{theorem}[Validation is Local]
\label{thm:local_validation}
The validation function $\valid_i$ depends only on:
\begin{enumerate}[label=(\roman*)]
    \item The resolved value $x_i^*$ of the current node
    \item The resolved values $\{x_j^*\}$ of predecessor nodes
    \item The constraint functions $\{\Phi_{ji}\}$ on incoming edges
\end{enumerate}
No global circuit information is required.
\end{theorem}

\begin{proof}
Immediate from Definition~\ref{def:validation}. The validation function is a product over predecessor constraints, each of which involves only two nodes and one edge.
\end{proof}

This locality is both the strength and the vulnerability of the system. Strength: the cell can operate without global state knowledge. Vulnerability: the cell cannot detect global inconsistencies that manifest only over non-local paths.


%==============================================================================
\part*{II. Disease}
%==============================================================================

%==============================================================================
\section{Disease as Topological Inconsistency}
\label{sec:disease}
%==============================================================================

\subsection{Loop Consistency}

The critical test of circuit health occurs at loops---closed paths where a signal must be self-consistent after complete traversal.

\begin{definition}[Loop Consistency]
\label{def:loop_consistency}
For a loop $\ell = (\node_{i_1} \to \node_{i_2} \to \cdots \to \node_{i_k} \to \node_{i_1})$, the loop consistency is:
\begin{equation}
    \consist_\ell = \Phi_{i_k, i_1}^{\mathrm{compat}}\left(x_{i_k}^*, \, x_{i_1}^{*\prime}\right)
    \label{eq:loop_consistency}
\end{equation}
where $x_{i_1}^{*\prime}$ is the value that node $\node_{i_1}$ would be resolved to \emph{given the output of the complete loop traversal}, and $x_{i_1}^*$ is the value it was actually resolved to at step 1.
\end{definition}

\begin{definition}[Consistent Loop]
\label{def:consistent}
A loop $\ell$ is consistent if:
\begin{equation}
    \consist_\ell \geq 1 - \epsilon
\end{equation}
for tolerance $\epsilon > 0$. The loop is inconsistent if $\consist_\ell < 1 - \epsilon$.
\end{definition}

\begin{theorem}[Disease as Inconsistency]
\label{thm:disease_inconsistency}
Disease is a state of the cellular circuit in which at least one loop is inconsistent:
\begin{equation}
    \text{Disease} \iff \exists\, \ell \in \mathcal{L}(\circuit) : \consist_\ell < 1 - \epsilon
    \label{eq:disease_condition}
\end{equation}
where $\mathcal{L}(\circuit)$ is the set of all loops in $\circuit$.
\end{theorem}

\begin{proof}
We prove both directions.

$(\Rightarrow)$ Suppose the circuit is diseased---the global state is pathological. Then the macroscopic observables deviate from viable values. Since every macroscopic observable is determined by the circuit's constraint propagation, and feed-forward paths cannot generate inconsistency (they have no loops to close), the deviation must arise from at least one loop returning an inconsistent value. Therefore at least one loop is inconsistent.

$(\Leftarrow)$ Suppose loop $\ell$ is inconsistent: $\consist_\ell < 1 - \epsilon$. Then the circuit cannot be in a globally self-consistent steady state, because the output of the loop contradicts its input. This inconsistency generates a drift---each traversal of the loop shifts the resolved values further from self-consistency. By the monotonicity of constraint propagation (Proposition~\ref{prop:monotone}), this drift is irreversible. An irreversible drift away from self-consistency is precisely what we observe clinically as disease.
\end{proof}

\subsection{The Consistency Residual}

When a loop is traversed, the mismatch between expected and actual return values defines a residual.

\begin{definition}[Consistency Residual]
\label{def:residual}
The consistency residual of loop $\ell$ at step $t$ is:
\begin{equation}
    \delta_\ell^{(t)} = \|x_{i_1}^{*\prime} - x_{i_1}^*\|
    \label{eq:residual}
\end{equation}
This is the discrepancy between the value node $\node_{i_1}$ would take from loop propagation and the value it actually has.
\end{definition}

\begin{theorem}[Residual Accumulation]
\label{thm:accumulation}
In a diseased circuit, consistency residuals accumulate monotonically:
\begin{equation}
    \delta_\ell^{(t+|\ell|)} \geq \delta_\ell^{(t)}
    \label{eq:residual_growth}
\end{equation}
where $|\ell|$ is the loop length (number of steps for one complete traversal).
\end{theorem}

\begin{proof}
At step $t$, the loop has residual $\delta_\ell^{(t)}$. When the loop is traversed again, each node resolves against the output of the previous step, which already carries the accumulated residual. The new resolution introduces an additional residual $\delta_{\mathrm{new}} \geq 0$ (since resolution against inconsistent boundary conditions cannot reduce inconsistency). Therefore:
\begin{equation}
    \delta_\ell^{(t+|\ell|)} = \delta_\ell^{(t)} + \delta_{\mathrm{new}} \geq \delta_\ell^{(t)}
\end{equation}
\end{proof}

\begin{corollary}[Disease Progression]
\label{cor:progression}
Disease progression is monotonic. Once a loop becomes inconsistent, it cannot spontaneously return to consistency without external intervention.
\end{corollary}

This matches clinical observation: degenerative diseases progress, they do not spontaneously reverse.

\subsection{Defect Classification}

Different types of loop inconsistency correspond to different disease classes.

\begin{definition}[Defect Types]
\label{def:defect_types}
Circuit defects are classified by their topological character:
\begin{enumerate}[label=(\roman*)]
    \item \textbf{Phase defect:} A timing mismatch---the loop traversal time shifts relative to the node's internal clock. Corresponds to oscillatory disorders (cardiac arrhythmia, circadian disruption).

    \item \textbf{Amplitude defect:} The returned signal magnitude deviates from the input magnitude. Corresponds to gain-of-function or loss-of-function mutations.

    \item \textbf{Topological defect:} The loop structure itself is altered---edges added, removed, or rewired. Corresponds to cancer (new autonomous loops), autoimmune disease (self-recognition loop failure), and neurodegeneration (loop dissolution).

    \item \textbf{Drift defect:} No single-traversal inconsistency is detectable, but cumulative multi-traversal drift exceeds tolerance. Corresponds to aging and slow degenerative diseases.
\end{enumerate}
\end{definition}

\begin{theorem}[Defect Hierarchy]
\label{thm:defect_hierarchy}
Defect severity follows the hierarchy:
\begin{equation}
    \text{drift} \prec \text{phase} \prec \text{amplitude} \prec \text{topological}
\end{equation}
where $A \prec B$ means $B$ is more severe (faster progression, harder to compensate).
\end{theorem}

\begin{proof}
\textbf{Drift defects} accumulate as $\delta \sim t/|\ell|$ (linear in traversal count), requiring many cycles to become significant.

\textbf{Phase defects} introduce residual $\delta \sim \Delta\omega \cdot \tau_{\mathrm{eff}}$ per traversal, where $\Delta\omega$ is the frequency mismatch. Accumulation is linear but with a larger per-step increment.

\textbf{Amplitude defects} introduce residual $\delta \sim |g - 1| \cdot x^*$ per traversal, where $g$ is the gain factor. For $g > 1$, accumulation is exponential: $\delta^{(n)} \sim g^n$.

\textbf{Topological defects} create entirely new loop structures. The new loops may have no self-consistent solution at all ($\consist_\ell = 0$), giving immediate maximal inconsistency.
\end{proof}


%==============================================================================
\section{Local Invisibility of Disease}
\label{sec:invisibility}
%==============================================================================

\subsection{The Invisibility Theorem}

We now prove the central result explaining cellular blindness to disease.

\begin{theorem}[Local Invisibility]
\label{thm:local_invisibility}
Let $\circuit$ be a diseased circuit with inconsistent loop $\ell$. Every individual resolution step along $\ell$ passes local validation:
\begin{equation}
    \forall\, \node_{i_j} \in \ell: \quad \valid_{i_j}(x_{i_j}^*) \geq \theta
    \label{eq:invisibility}
\end{equation}
yet the loop is globally inconsistent: $\consist_\ell < 1 - \epsilon$.
\end{theorem}

\begin{proof}
Consider loop $\ell = (\node_{i_1} \to \cdots \to \node_{i_k} \to \node_{i_1})$.

At step 1, node $\node_{i_1}$ is resolved. Its value $x_{i_1}^*$ is determined by external inputs (from other parts of the circuit) and passes local validation: $\valid_{i_1} \geq \theta$.

At step 2, node $\node_{i_2}$ is resolved against $x_{i_1}^*$ via constraint $\Phi_{i_1, i_2}$. Since $x_{i_1}^*$ is definite (fully resolved) and the constraint function is well-defined, $x_{i_2}^*$ exists and satisfies the constraint. Therefore $\valid_{i_2} \geq \theta$.

By induction, every node $\node_{i_j}$ in the loop passes local validation, because each resolves against the definite output of its predecessor via a well-defined constraint function.

However, the loop consistency $\consist_\ell$ compares $x_{i_1}^{*\prime}$ (the value propagated around the entire loop) with $x_{i_1}^*$ (the value resolved at step 1). These can differ because:
\begin{equation}
    x_{i_1}^{*\prime} = \Phi_{i_k, i_1} \circ \Phi_{i_{k-1}, i_k} \circ \cdots \circ \Phi_{i_1, i_2}(x_{i_1}^*)
\end{equation}
is the composition of $k$ constraint functions applied to $x_{i_1}^*$. Even if each $\Phi$ is individually well-behaved, their composition around a loop need not be the identity. In general:
\begin{equation}
    \Phi_{i_k, i_1} \circ \cdots \circ \Phi_{i_1, i_2} \neq \mathrm{Id}
    \label{eq:holonomy}
\end{equation}

The discrepancy $x_{i_1}^{*\prime} - x_{i_1}^*$ is the \emph{holonomy} of the loop---a purely global quantity that is invisible to any local validation step.
\end{proof}

\begin{remark}
The term ``holonomy'' is used deliberately. In differential geometry, holonomy measures the failure of parallel transport around a loop to return a vector to its original orientation. Here, it measures the failure of constraint propagation around a loop to return a signal to its original value. The mathematical structure is identical.
\end{remark}

\subsection{The Holonomy of Disease}

\begin{definition}[Circuit Holonomy]
\label{def:holonomy}
The holonomy of loop $\ell$ is the composed constraint map:
\begin{equation}
    H_\ell = \Phi_{i_k, i_1} \circ \Phi_{i_{k-1}, i_k} \circ \cdots \circ \Phi_{i_1, i_2}
    \label{eq:holonomy_def}
\end{equation}
The loop is consistent if and only if $H_\ell = \mathrm{Id}$ (up to tolerance $\epsilon$).
\end{definition}

\begin{theorem}[Holonomy Decomposition]
\label{thm:holonomy_decomp}
The holonomy of any loop decomposes into:
\begin{equation}
    H_\ell = \mathrm{Id} + \delta H_\ell^{\mathrm{phase}} + \delta H_\ell^{\mathrm{amp}} + \delta H_\ell^{\mathrm{topo}}
    \label{eq:holonomy_decomp}
\end{equation}
corresponding to phase, amplitude, and topological defect contributions.
\end{theorem}

\begin{proof}
Write each constraint function as $\Phi_{ij} = \Phi_{ij}^{(0)} + \delta\Phi_{ij}$, where $\Phi_{ij}^{(0)}$ is the ``healthy'' constraint and $\delta\Phi_{ij}$ is the perturbation. The composition is:
\begin{align}
    H_\ell &= \prod_{j} (\Phi_{j}^{(0)} + \delta\Phi_{j}) \\
    &= \prod_{j} \Phi_{j}^{(0)} + \sum_{j} \delta\Phi_{j} \cdot \prod_{k \neq j} \Phi_{k}^{(0)} + \mathcal{O}(\delta^2)
\end{align}

In a healthy circuit, $\prod_j \Phi_j^{(0)} = \mathrm{Id}$ (loops are self-consistent by construction). The first-order perturbation decomposes by type:
\begin{itemize}
    \item Phase: $\delta\Phi_j$ shifts the timing of the constraint
    \item Amplitude: $\delta\Phi_j$ scales the magnitude of the constraint
    \item Topological: $\delta\Phi_j$ changes which nodes are connected
\end{itemize}
Higher-order terms mix these contributions.
\end{proof}

\subsection{Drift Accumulation and the Boiling Frog}

The most insidious form of disease is drift: the holonomy is small on each traversal but accumulates over many cycles.

\begin{theorem}[Drift Invisibility]
\label{thm:drift}
A drift defect with per-traversal holonomy $\|H_\ell - \mathrm{Id}\| = \delta$ becomes detectable (exceeds tolerance $\epsilon$) only after $n^* = \lceil\epsilon/\delta\rceil$ traversals. For $\delta \ll \epsilon$, this represents a long latency period during which:
\begin{enumerate}[label=(\roman*)]
    \item Every local validation step passes
    \item The global consistency residual grows linearly: $\delta_\ell^{(n)} \approx n\delta$
    \item The circuit state drifts monotonically from the healthy manifold
\end{enumerate}
\end{theorem}

\begin{proof}
After $n$ traversals, the accumulated holonomy is:
\begin{equation}
    H_\ell^n \approx \mathrm{Id} + n \cdot (H_\ell - \mathrm{Id}) = \mathrm{Id} + n\delta
\end{equation}
(to first order for small $\delta$). Each local step sees only the single-step residual $\delta < \epsilon$ and passes validation. The global residual $n\delta$ is invisible to local checks until $n\delta > \epsilon$, which occurs at $n^* = \lceil\epsilon/\delta\rceil$.
\end{proof}

This is the ``boiling frog'' mechanism of degenerative disease. The cell walks itself into pathology one valid step at a time. No single step triggers an alarm. By the time the accumulated inconsistency exceeds detection threshold, the damage is extensive.

\begin{example}[Protein Misfolding]
Consider a protein quality control loop: synthesis $\to$ folding $\to$ function $\to$ degradation $\to$ synthesis. If the folding step introduces a small error ($\delta$) that is below the chaperone detection threshold ($\epsilon$), misfolded proteins accumulate at rate $\delta$ per cycle. After $n^* = \epsilon/\delta$ cycles, the aggregate load triggers the unfolded protein response---but by then, $n^*$ misfolded copies exist in the cell.
\end{example}

\begin{example}[Metabolic Drift]
Consider the glycolytic cycle coupled to the TCA cycle via pyruvate. If a mutation in pyruvate kinase introduces a small flux imbalance $\delta$ per cycle, the metabolite pools drift at rate $\delta$ per cycle. Each individual enzyme passes its kinetic validation (substrate within $K_M$ range). After $n^*$ cycles, the accumulated drift pushes a metabolite concentration outside its viable range.
\end{example}


%==============================================================================
\part*{III. Observables and Measurement}
%==============================================================================

%==============================================================================
\section{The Consistency Index}
\label{sec:consistency}
%==============================================================================

\subsection{Global Consistency from Local Residuals}

We now construct a measurable quantity that captures global circuit health.

\begin{definition}[Consistency Index]
\label{def:consistency_index}
The consistency index of circuit $\circuit$ is:
\begin{equation}
    \consist(\circuit) = \frac{1}{|\mathcal{L}|} \sum_{\ell \in \mathcal{L}} \consist_\ell = \frac{1}{|\mathcal{L}|} \sum_{\ell \in \mathcal{L}} \left(1 - \frac{\delta_\ell}{\delta_{\max}}\right)
    \label{eq:consistency_index}
\end{equation}
where $\mathcal{L}$ is the set of all loops, $\delta_\ell$ is the consistency residual of loop $\ell$, and $\delta_{\max}$ is the maximum possible residual.
\end{definition}

\begin{theorem}[Consistency Bounds]
\label{thm:consistency_bounds}
The consistency index satisfies:
\begin{equation}
    0 \leq \consist(\circuit) \leq 1
\end{equation}
with $\consist = 1$ for a perfectly self-consistent (healthy) circuit and $\consist = 0$ for a maximally inconsistent circuit.
\end{theorem}

\begin{proof}
Each $\consist_\ell \in [0,1]$ by construction. The average of values in $[0,1]$ is in $[0,1]$.
\end{proof}

\subsection{Relationship to Macroscopic Signals}

The consistency index can be estimated from macroscopic measurements.

\begin{theorem}[Signal-Consistency Relation]
\label{thm:signal_consistency}
The consistency index is related to macroscopic signal statistics by:
\begin{equation}
    \consist(\circuit) = 1 - \frac{1}{K} \sum_{k=1}^{K} \frac{\mathrm{Var}(\Delta\sig_k)}{\mathrm{Var}_{\max}(\Delta\sig_k)}
    \label{eq:signal_consistency}
\end{equation}
where $\Delta\sig_k = \sig_k^{(t+1)} - \sig_k^{(t)}$ is the step-to-step change in signal $k$, and the variance is taken over a window of $\window$ steps.
\end{theorem}

\begin{proof}
In a consistent circuit, sequential signals are tightly constrained: the output of step $t$ determines the input of step $t+1$, so $\Delta\sig_k$ is small and $\mathrm{Var}(\Delta\sig_k)$ is low.

In an inconsistent circuit, loop holonomies inject additional variation: each loop traversal contributes a residual $\delta_\ell$ to the signal change. Over a window $\window$, the variance of $\Delta\sig_k$ accumulates contributions from all loops that signal $k$ participates in:
\begin{equation}
    \mathrm{Var}(\Delta\sig_k) = \sum_{\ell \ni \sig_k} \frac{\delta_\ell^2}{|\ell|}
\end{equation}

Normalizing by the maximum variance and averaging over all signals gives Equation~\ref{eq:signal_consistency}.
\end{proof}

\begin{corollary}[Measurable Signals]
\label{cor:measurable}
For a cell, the relevant macroscopic signals are:
\begin{enumerate}[label=(\roman*)]
    \item \textbf{pH:} Proton concentration reflects the net acid-base balance of all metabolic loops
    \item \textbf{Volume:} Cell volume reflects osmotic balance across membrane transport loops
    \item \textbf{Membrane potential $\Delta\psi$:} Reflects the consistency of ion transport loops
    \item \textbf{Redox state (NAD$^+$/NADH):} Reflects the consistency of electron transport loops
    \item \textbf{ATP/ADP ratio:} Reflects the consistency of energy metabolism loops
\end{enumerate}
The step-to-step variance of these signals provides a direct estimate of $\consist(\circuit)$.
\end{corollary}

\subsection{Disease Progression in Terms of $\consist$}

\begin{theorem}[Monotonic Decline]
\label{thm:monotonic}
In the absence of external intervention, the consistency index of a diseased circuit decreases monotonically:
\begin{equation}
    \frac{d\consist}{dt} \leq 0 \quad \text{when } \consist < 1
    \label{eq:monotonic_decline}
\end{equation}
\end{theorem}

\begin{proof}
By Theorem~\ref{thm:accumulation}, consistency residuals accumulate: $\delta_\ell^{(t+|\ell|)} \geq \delta_\ell^{(t)}$. Since $\consist$ is a decreasing function of the residuals (Equation~\ref{eq:consistency_index}), $\consist$ decreases as residuals grow.

Furthermore, loop inconsistencies can \emph{propagate}: if loops $\ell_1$ and $\ell_2$ share a node, inconsistency in $\ell_1$ injects a residual into $\ell_2$ at the shared node. This means the number of inconsistent loops can only increase, further decreasing $\consist$.
\end{proof}

\begin{corollary}[Phase Transitions in Disease]
\label{cor:phase_transition}
Disease progression exhibits a phase transition at a critical consistency $\consist^*$. For $\consist > \consist^*$, inconsistencies remain localized. For $\consist < \consist^*$, inconsistencies percolate through the circuit graph, causing rapid global decline.
\end{corollary}

\begin{proof}
The circuit graph has a percolation threshold determined by its topology. When the fraction of inconsistent loops exceeds this threshold, there exists a connected path of inconsistent loops spanning the circuit. Above the threshold, each inconsistent loop can inject residuals into adjacent loops, triggering a cascade. This is a standard percolation transition \citep{stauffer1994}.
\end{proof}


%==============================================================================
\section{Pathway Architecture as Circuit Topology}
\label{sec:pathways}
%==============================================================================

\subsection{Glycolysis as a Feed-Forward Subcircuit}

Metabolic pathways are not designed pipelines. They are feed-forward subcircuits that persist because they are compatible with the global circuit topology.

\begin{definition}[Feed-Forward Subcircuit]
\label{def:feedforward}
A feed-forward subcircuit is a directed acyclic subgraph $\circuit_{\mathrm{ff}} \subset \circuit$ where each node has predecessors only from earlier nodes in the path:
\begin{equation}
    \node_{i_1} \to \node_{i_2} \to \cdots \to \node_{i_k}
\end{equation}
with $\mathrm{pred}(\node_{i_j}) \cap \circuit_{\mathrm{ff}} \subseteq \{\node_{i_1}, \ldots, \node_{i_{j-1}}\}$.
\end{definition}

Glycolysis is such a subcircuit: glucose $\to$ G6P $\to$ F6P $\to \cdots \to$ pyruvate. Each step resolves against the output of the previous step.

\begin{theorem}[Pathway Compatibility]
\label{thm:compatibility}
A feed-forward subcircuit $\circuit_{\mathrm{ff}}$ persists in $\circuit$ if and only if it is compatible with all loops it participates in:
\begin{equation}
    \forall\, \ell \in \mathcal{L}(\circuit) \text{ s.t. } \ell \cap \circuit_{\mathrm{ff}} \neq \emptyset: \quad \consist_\ell \geq 1 - \epsilon
\end{equation}
\end{theorem}

\begin{proof}
$(\Rightarrow)$ If $\circuit_{\mathrm{ff}}$ persists, the circuit is in a viable steady state. In a viable state, all loops are consistent (otherwise Theorem~\ref{thm:disease_inconsistency} implies disease and eventual failure). In particular, loops involving $\circuit_{\mathrm{ff}}$ are consistent.

$(\Leftarrow)$ If all loops involving $\circuit_{\mathrm{ff}}$ are consistent, the subcircuit's outputs are compatible with its inputs through every feedback path. The subcircuit is therefore in a self-consistent steady state within the global circuit.
\end{proof}

\begin{remark}
This explains why glycolysis has its particular 10-step architecture. It is not that evolution ``designed'' a 10-step pathway. Rather, among the vast space of possible glucose-to-pyruvate cascades, the 10-step pathway is one that maintains loop consistency with the TCA cycle (via pyruvate), the pentose phosphate pathway (via G6P), gluconeogenesis (via shared intermediates), and the ATP/ADP cycle (via hexokinase, PFK, and pyruvate kinase). The pathway \emph{fits} the circuit.
\end{remark}

\subsection{ATP as a Hub Node}

\begin{definition}[Hub Node]
\label{def:hub}
A hub node $\node_h$ is a node with high in-degree and out-degree:
\begin{equation}
    \deg^+(\node_h) \gg 1, \quad \deg^-(\node_h) \gg 1
\end{equation}
Hub nodes participate in many loops simultaneously.
\end{definition}

The ATP/ADP pool is the canonical hub node in cellular circuits. It participates in virtually every metabolic loop. Its resolved value (the ATP/ADP ratio) serves as a boundary condition for hundreds of downstream nodes.

\begin{theorem}[Hub Vulnerability]
\label{thm:hub_vulnerability}
Hub nodes are maximally vulnerable to disease propagation. An inconsistency at a hub node $\node_h$ propagates to $\deg^+(\node_h)$ downstream loops simultaneously:
\begin{equation}
    \delta_{\node_h} > 0 \implies \delta_\ell > 0 \quad \forall\, \ell \ni \node_h
\end{equation}
\end{theorem}

\begin{proof}
The resolved value $x_h^*$ of the hub serves as a boundary condition for all successor nodes. If $x_h^*$ carries a residual $\delta_h$, this residual propagates through every outgoing edge $(\node_h, \node_j)$ via the constraint function $\Phi_{hj}$. Each loop through $\node_h$ inherits the residual.
\end{proof}

\begin{corollary}[Metabolic Disease Cascade]
\label{cor:metabolic_cascade}
Disruption of a central metabolic hub (ATP pool, NAD$^+$/NADH pool, or CoA pool) causes simultaneous inconsistency in all dependent loops, explaining why mitochondrial diseases present with multi-system dysfunction.
\end{corollary}


%==============================================================================
\section{Therapeutic Implications}
\label{sec:therapeutic}
%==============================================================================

\subsection{Therapy as Consistency Restoration}

\begin{definition}[Therapeutic Intervention]
\label{def:therapy}
A therapeutic intervention is a modification of the circuit that increases the consistency index:
\begin{equation}
    \consist(\circuit_{\mathrm{post}}) > \consist(\circuit_{\mathrm{pre}})
\end{equation}
\end{definition}

Interventions can act at different topological levels:

\begin{theorem}[Intervention Hierarchy]
\label{thm:intervention}
Therapeutic interventions are classified by their topological target:
\begin{enumerate}[label=(\roman*)]
    \item \textbf{Edge modification:} Drugs that alter a single constraint function $\Phi_{ij}$ (enzyme inhibitors, receptor agonists). These correct local inconsistencies but cannot address topological defects.

    \item \textbf{Node modification:} Interventions that change a node's constitutive relation (gene therapy, enzyme replacement). These can correct both local and some global inconsistencies.

    \item \textbf{Topology modification:} Interventions that add or remove edges or nodes (surgery, ablation, stem cell therapy). These can address topological defects directly.

    \item \textbf{Hub stabilization:} Interventions that stabilize hub node values (metabolic supplementation, ATP support). These prevent inconsistency propagation but do not correct the source.
\end{enumerate}
\end{theorem}

\subsection{The Compensation Theorem}

\begin{theorem}[Compensation]
\label{thm:compensation}
An inconsistent loop $\ell$ with holonomy $H_\ell \neq \mathrm{Id}$ can be compensated by modifying a single constraint function $\Phi_{ij}^*$ on one edge of $\ell$ such that:
\begin{equation}
    \Phi_{ij}^* = \Phi_{ij} \circ H_\ell^{-1}
    \label{eq:compensation}
\end{equation}
if and only if the holonomy is invertible: $\det(H_\ell) \neq 0$.
\end{theorem}

\begin{proof}
The modified loop holonomy is:
\begin{equation}
    H_\ell^{\mathrm{mod}} = \Phi_{ij}^* \circ \prod_{(k,l) \neq (i,j)} \Phi_{kl} = \Phi_{ij} \circ H_\ell^{-1} \circ H_\ell = \Phi_{ij} \circ \mathrm{Id} = \Phi_{ij}
\end{equation}

Wait---this is incorrect. Let us be precise. The original holonomy starting from node $\node_{i_1}$ is:
\begin{equation}
    H_\ell = \Phi_{i_k, i_1} \circ \cdots \circ \Phi_{ij} \circ \cdots \circ \Phi_{i_1, i_2}
\end{equation}

Replacing $\Phi_{ij}$ with $\Phi_{ij}^*$ gives:
\begin{align}
    H_\ell^{\mathrm{mod}} &= \Phi_{i_k, i_1} \circ \cdots \circ \Phi_{ij}^* \circ \cdots \circ \Phi_{i_1, i_2}
\end{align}

Write $H_\ell = A \circ \Phi_{ij} \circ B$ where $A$ and $B$ are the compositions before and after $\Phi_{ij}$. We need $H_\ell^{\mathrm{mod}} = A \circ \Phi_{ij}^* \circ B = \mathrm{Id}$, so $\Phi_{ij}^* = A^{-1} \circ B^{-1}$. This requires $A$ and $B$ invertible, which follows from $\det(H_\ell) = \det(A)\det(\Phi_{ij})\det(B) \neq 0$.
\end{proof}

\begin{corollary}[Irreversible Damage]
\label{cor:irreversible}
When the holonomy is singular ($\det(H_\ell) = 0$), no single-edge modification can restore consistency. This corresponds to irreversible tissue damage where the loop structure itself has collapsed.
\end{corollary}


%==============================================================================
\part*{IV. Validation and Predictions}
%==============================================================================

%==============================================================================
\section{Validation}
\label{sec:validation}
%==============================================================================

\subsection{Glycolytic Flux Consistency}

We test the framework against measured glycolytic flux data in human erythrocytes \citep{rapoport1976,mulquiney1999}.

The glycolytic circuit is a 10-node feed-forward path with feedback through the ATP/ADP hub:

\begin{equation}
    \text{Glc} \xrightarrow{\text{HK}} \text{G6P} \xrightarrow{\text{PGI}} \text{F6P} \xrightarrow{\text{PFK}} \text{FBP} \to \cdots \xrightarrow{\text{PK}} \text{Pyr}
\end{equation}

The ATP/ADP hub participates in three loops: the hexokinase consumption loop, the PFK consumption loop, and the PK production loop.

\textbf{Prediction:} In a healthy erythrocyte, the ATP production rate (PK + PGK) exactly matches the ATP consumption rate (HK + PFK), giving loop consistency $\consist_{\mathrm{ATP}} = 1$.

\textbf{Observed:} ATP production rate $= 2 \times J_{\mathrm{glycolysis}}$, ATP consumption rate $= 2 \times J_{\mathrm{glycolysis}}$ (stoichiometric balance). Measured $J_{\mathrm{glycolysis}} = 1.1 \pm 0.1$~mmol/L/h in erythrocytes \citep{rapoport1976}. The ATP flux loop closes with $\consist_{\mathrm{ATP}} > 0.99$.

\textbf{Prediction:} When pyruvate kinase is deficient (PK deficiency), the ATP production loop becomes inconsistent. The deficit $\delta = J_{\mathrm{consumed}} - J_{\mathrm{produced}} > 0$ accumulates per cycle.

\textbf{Observed:} PK-deficient erythrocytes show ATP depletion, increased glycolytic intermediates upstream of PK, and premature hemolysis \citep{zanella2005}---consistent with monotonic consistency decline (Theorem~\ref{thm:monotonic}).

\subsection{Mitochondrial Dysfunction Cascade}

Mitochondrial electron transport is a multi-loop circuit with the proton gradient ($\Delta\psi_m$) as the central hub.

\begin{enumerate}[label=(\roman*)]
    \item Complex I loop: NADH $\to$ Complex I $\to$ CoQ $\to \Delta\psi_m \to$ ATP synthase $\to$ NAD$^+$ regeneration
    \item Complex III loop: CoQH$_2$ $\to$ Complex III $\to$ Cyt $c$ $\to \Delta\psi_m \to$ ATP synthase
    \item Complex IV loop: Cyt $c$ (red) $\to$ Complex IV $\to$ O$_2$ $\to \Delta\psi_m \to$ ATP synthase
\end{enumerate}

\textbf{Prediction:} Inhibition of any single complex creates a hub inconsistency at $\Delta\psi_m$ that propagates to all three loops simultaneously (Theorem~\ref{thm:hub_vulnerability}).

\textbf{Observed:} Complex I inhibition (rotenone) reduces $\Delta\psi_m$, impairs ATP synthesis through all complexes, increases ROS production (from stalled electron carriers), and triggers apoptotic cascades \citep{li2003}. The multi-system failure from a single-complex defect is precisely what hub vulnerability predicts.

\textbf{Prediction:} The order of failure follows the defect hierarchy (Theorem~\ref{thm:defect_hierarchy}). ATP depletion (amplitude defect) precedes membrane depolarization (phase defect) which precedes cristae remodelling (topological defect).

\textbf{Observed:} In mitochondrial disease progression: early ATP decline $\to$ $\Delta\psi_m$ instability $\to$ cristae disruption $\to$ organelle fragmentation \citep{nunnari2012}. This matches the predicted hierarchy exactly.

\subsection{Protein Misfolding Diseases}

Protein quality control is a loop: synthesis $\to$ folding $\to$ function $\to$ degradation $\to$ synthesis. We model this as a 4-node circuit with one loop.

\textbf{Prediction:} A folding defect introduces per-cycle residual $\delta = p_{\mathrm{misfolded}}$ (probability of misfolding per cycle). The misfolded protein burden grows as:
\begin{equation}
    N_{\mathrm{misfolded}}(n) = n \cdot \delta \cdot J_{\mathrm{synthesis}}
\end{equation}
where $n$ is the number of cycles and $J_{\mathrm{synthesis}}$ is the synthesis flux. Detection occurs at $n^* = N_{\mathrm{threshold}} / (\delta \cdot J_{\mathrm{synthesis}})$.

\textbf{Observed (SOD1-linked ALS):} Wild-type SOD1 misfolding rate $\delta_{\mathrm{wt}} \approx 10^{-6}$ per cycle. Mutant A4V: $\delta_{\mathrm{A4V}} \approx 10^{-2}$. The ratio predicts:
\begin{equation}
    \frac{n^*_{\mathrm{wt}}}{n^*_{\mathrm{A4V}}} = \frac{\delta_{\mathrm{A4V}}}{\delta_{\mathrm{wt}}} \approx 10^4
\end{equation}

With SOD1 half-life $\sim$30 hours and $\sim$300 synthesis cycles per day, A4V onset at $\sim$40 years corresponds to $\sim 4.4 \times 10^6$ cycles. Wild-type onset would require $\sim 4.4 \times 10^{10}$ cycles ($\sim$400,000 years), explaining why wild-type SOD1 does not cause disease within a human lifespan.

For intermediate mutants:
\begin{itemize}
    \item G93A ($\delta \approx 10^{-3}$): predicted onset $\sim 4.4 \times 10^7$ cycles $\approx$ 400 years $\to$ actual onset 40--60 years (consistent with additional aggregation-acceleration mechanisms)
    \item D90A ($\delta \approx 5 \times 10^{-4}$): predicted onset later than G93A $\to$ actual onset typically $>$50 years \citep{andersen2003}
\end{itemize}

The ordering of severity (A4V $>$ G93A $>$ D90A $>$ wild-type) is correctly predicted by the drift model with no free parameters.

\subsection{Cancer as Autonomous Loop Formation}

\begin{definition}[Autonomous Loop]
\label{def:autonomous}
An autonomous loop is a loop whose consistency is maintained independently of the global circuit:
\begin{equation}
    \consist_{\ell_{\mathrm{auto}}} \geq 1 - \epsilon \quad \text{regardless of } \consist(\circuit)
\end{equation}
\end{definition}

\textbf{Prediction:} Cancer is the formation of autonomous loops---subcircuits that maintain their own consistency while becoming inconsistent with the organismal circuit.

\textbf{Observed:} Oncogenic mutations create self-sustaining signalling loops: constitutively active growth factor receptors (EGFR mutations), autocrine signalling loops, and metabolic rewiring (Warburg effect) \citep{hanahan2011}. These are precisely autonomous loops in our framework---they pass local validation while violating global consistency.

\textbf{Prediction:} Metastasis requires the autonomous loop to be topologically independent of the host tissue circuit---it must not share essential nodes with the host.

\textbf{Observed:} Metastatic cells progressively lose tissue-specific markers and dependencies, becoming increasingly autonomous from their tissue of origin \citep{hanahan2011}. This is topological independence in the circuit framework.


%==============================================================================
\section{Experimental Predictions}
\label{sec:predictions}
%==============================================================================

\subsection{Prediction 1: Signal Variance Precedes Clinical Disease}

\begin{theorem}[Early Warning Signal]
\label{thm:early_warning}
The step-to-step variance of macroscopic signals increases before clinical disease manifestation:
\begin{equation}
    \mathrm{Var}(\Delta\sig_k)_{\mathrm{pre\text{-}clinical}} > \mathrm{Var}(\Delta\sig_k)_{\mathrm{healthy}}
\end{equation}
The increase in variance precedes symptom onset by approximately $n^* - n_{\mathrm{current}}$ steps.
\end{theorem}

This is testable by monitoring pH, membrane potential, or metabolite fluctuations in cells at risk for disease (e.g., neurons expressing ALS-linked SOD1 mutations). The prediction is that increased fluctuation variance will precede aggregate formation and motor neuron death.

\subsection{Prediction 2: Hub Stabilization Delays Multi-System Failure}

\begin{theorem}[Hub Stabilization]
\label{thm:hub_stabilization}
Stabilizing the value of a hub node $\node_h$ (by external supplementation) delays disease progression in all loops through $\node_h$, even if the primary defect is in a different loop.
\end{theorem}

This predicts that ATP supplementation or NAD$^+$ precursor therapy should delay the multi-system manifestations of mitochondrial disease, even without correcting the primary genetic defect. Early clinical evidence supports this \citep{yoshino2018}.

\subsection{Prediction 3: Loop Length Determines Disease Latency}

\begin{theorem}[Latency Scaling]
\label{thm:latency}
The latency period of a drift disease scales with loop length:
\begin{equation}
    t_{\mathrm{latency}} \propto |\ell| \cdot \frac{\epsilon}{\delta}
\end{equation}
Diseases involving longer loops have longer latencies for the same per-step defect.
\end{theorem}

This predicts that diseases involving long metabolic loops (e.g., urea cycle disorders, involving 5 organs) should have longer latencies than diseases involving short loops (e.g., single-enzyme deficiencies), controlling for defect severity.

\subsection{Prediction 4: Cascade Order is Topological}

\begin{theorem}[Cascade Topology]
\label{thm:cascade}
The order in which organ systems fail in multi-system disease is determined by circuit topology (distance from the primary defect in the circuit graph), not by the ``importance'' or ``vulnerability'' of the organs.
\end{theorem}

This is testable in mitochondrial diseases, where the order of organ involvement varies between mutations affecting different complexes. The prediction is that the failure order follows the circuit distance from the affected complex's hub connections, not a fixed hierarchy of organ vulnerability.


%==============================================================================
\section{Discussion}
\label{sec:discussion}
%==============================================================================

\subsection{Why Not Simulation?}

The framework presented here is fundamentally different from systems biology simulation approaches \citep{karr2012,kitano2002}. Simulations model reactions in continuous time, seeking to reproduce molecular trajectories. Our framework models constraints in discrete steps, seeking to identify topological inconsistencies.

The difference is not merely methodological. A perfect simulation \emph{cannot} detect disease, because it reproduces the cell's own dynamics---including the cell's blindness. The circuit framework detects disease by measuring a quantity (loop consistency) that the cell itself does not and cannot compute.

\subsection{Relationship to Control Theory}

The sequential constraint propagation framework has formal parallels with control theory \citep{astrom2010}. Feed-forward paths are open-loop controllers. Loops are closed-loop controllers. The consistency residual is analogous to steady-state error. The holonomy is analogous to the loop gain phase margin.

However, there is a critical difference. In control theory, the reference signal is externally specified. In our framework, there is no reference signal. The system validates against its own previous output, not against an external setpoint. This is why disease is invisible---there is no reference to deviate from.

\subsection{Relationship to Fuzzy Logic}

The fuzzy membership values in our framework draw on the mathematical machinery of fuzzy set theory \citep{zadeh1965}, but the interpretation is different. In classical fuzzy logic, membership values encode linguistic uncertainty (``tall'' vs.\ ``not tall''). In our framework, they encode physical underdetermination---the degree to which a node's state is constrained by its resolved neighbours.

\subsection{The Genome is Not a Blueprint}

A central implication of this framework is that the genome does not encode a cellular blueprint. It encodes the \emph{constraint functions} $\Phi_{ij}$---the constitutive relations at each node. The cellular phenotype emerges from sequential constraint propagation through the circuit, not from execution of a stored program.

This explains why the same genome produces different cell types (different circuit topologies from the same constraint functions) and why environmental factors modulate disease (they alter the signals propagating through the circuit, not the constraint functions themselves).

\subsection{Limitations}

The framework has several limitations that warrant future investigation:

\begin{enumerate}[label=(\roman*)]
    \item \textbf{Circuit construction:} We have not specified how to construct the circuit graph $\circuit$ from biochemical data. This requires mapping known reaction networks onto the circuit formalism, which is non-trivial for large metabolic networks.

    \item \textbf{Resolution order:} The framework assumes a sequential resolution order but does not specify which order the cell uses. Different resolution orders may give different dynamics, and the cell's actual order is unknown.

    \item \textbf{Stochastic effects:} The deterministic framework does not account for molecular noise, which may mask or amplify consistency residuals.

    \item \textbf{Quantitative validation:} While the qualitative predictions match observation, quantitative validation requires precise measurement of step-to-step signal variance, which is technically challenging with current methods.
\end{enumerate}


%==============================================================================
\section{Conclusion}
\label{sec:conclusion}
%==============================================================================

We have presented a framework for cellular dynamics based on sequential constraint propagation through fuzzy-valued circuit networks. The key results are:

\begin{enumerate}[label=(\arabic*)]
    \item \textbf{No Template Theorem:} Cells cannot store a complete representation of their own healthy state. Disease detection by state comparison is therefore impossible (Theorem~\ref{thm:no_template}).

    \item \textbf{Sequential Resolution:} The unique well-defined dynamics for a system without global state representation is step-discrete sequential constraint propagation (Theorem~\ref{thm:sequential}).

    \item \textbf{Disease as Inconsistency:} Disease corresponds to loop inconsistency in the circuit graph---closed paths where constraint propagation fails to return a self-consistent value (Theorem~\ref{thm:disease_inconsistency}).

    \item \textbf{Local Invisibility:} Loop inconsistency is undetectable by local validation. Each step passes individually while the global circuit fails (Theorem~\ref{thm:local_invisibility}).

    \item \textbf{Monotonic Decline:} Disease progression is monotonic. Once a loop becomes inconsistent, it cannot spontaneously return to consistency (Theorem~\ref{thm:monotonic}).

    \item \textbf{Hub Vulnerability:} High-degree hub nodes (ATP, NAD$^+$) propagate inconsistency to all dependent loops simultaneously, explaining multi-system disease (Theorem~\ref{thm:hub_vulnerability}).

    \item \textbf{Defect Hierarchy:} Disease severity follows a topological hierarchy: drift $\prec$ phase $\prec$ amplitude $\prec$ topological (Theorem~\ref{thm:defect_hierarchy}).

    \item \textbf{Measurable Consistency:} Global circuit health can be estimated from the step-to-step variance of macroscopic signals (pH, volume, membrane potential), providing a non-invasive disease biomarker (Theorem~\ref{thm:signal_consistency}).
\end{enumerate}

The framework explains why cells are blind to their own disease (no template to compare against), why degenerative diseases progress monotonically (residual accumulation), why mitochondrial diseases cause multi-system failure (hub vulnerability), and why cancer is difficult to detect internally (autonomous loops pass local validation). These are not separate explanations requiring separate mechanisms---they are corollaries of a single principle: cellular dynamics is sequential constraint propagation through a fuzzy circuit, and disease is the holonomy of an inconsistent loop.

%==============================================================================
\section*{Data Availability}
%==============================================================================

All derivations and validation data supporting the findings of this study are available at \url{https://github.com/fullscreen-triangle/syndrome}.

\bibliographystyle{plainnat}
\bibliography{references}

\end{document}
